\documentclass[a4paper]{article}

\usepackage{graphicx}
\usepackage{amsmath,amssymb}
\usepackage{minted}
\usepackage{multirow}
\usepackage{float}
\usepackage{algorithm}
\usepackage{algpseudocode}
\usepackage{todonotes}
\usepackage{forest}
\usepackage{xstring}
\usepackage{xcolor}
\usepackage[most]{tcolorbox}
\usepackage{xparse}
\usepackage{natbib}

\usetikzlibrary{graphs,graphdrawing}
\usegdlibrary{layered}

% for external graphviz
\input{/home/donkarlo/Dropbox/repo/nd_language_ai_project/src/nd_language_ai/formal/computer/programming/tex/latex/code_clips/external_graphviz.tex}

% for tags
\input{/home/donkarlo/Dropbox/repo/nd_language_ai_project/src/nd_language_ai/formal/computer/programming/tex/latex/parts/control_sequence/kind/semtag/semtag.tex}

% for links
\input{/home/donkarlo/Dropbox/repo/nd_language_ai_project/src/nd_language_ai/formal/computer/programming/tex/latex/code_clips/seehere.tex}

\title{Navigator Clipboard, Pinned Ctrl+W, and Fresh Graphic Bridge Ownership}
\date{2026-09-06}
\author{Mohammad Rahmani}

\begin{document}
   \maketitle

   \section{Bugs}
   Three previously reported behaviors were still not effective at runtime:
   the later WorkspaceInteractionController replaced the multi-path Navigator
   menu, pinned tabs rejected explicit Ctrl+W close requests, and a stale
   graphic bridge on the default port could prevent a newly launched Mind
   Mirror process from running the current bridge code.

   \section{Fix}
   \begin{itemize}
      \item A dedicated NavigatorMultiPathOwnershipController becomes the final
      runtime owner of Ctrl+C, Ctrl+X, Ctrl+V, ExtendedSelection, and the
      multi-path right-click menu after all startup controllers are installed.
      \item Explicit Ctrl+W temporarily removes pin protection, uses the normal
      unsaved-changes close workflow, and restores the pin if the close is
      cancelled.
      \item Each Mind Mirror application session now starts one bridge process
      from the current project source. The bridge entry point already selects
      collision-safe ports, so a stale process on port 8766 no longer prevents
      current bridge code from taking over the iPad reverse connection.
   \end{itemize}

   \bibliographystyle{plainnat}
   \bibliography{/home/donkarlo/Dropbox/repo/nd_shared_research/bibliography/bibliography.bib}
\end{document}
